% GENERATED FILE. Edit canonical lessons, then run npm run generate:books.
% canonical-source-hash: fnv1a64:bb7bfe76444c94fc
% canonical-lessons: KA-C226-kale, KA-C226-atte, KA-C226-sose, KA-C226-aliya, KA-C226-attige

\chapter{Art, Aunt, Daughter-in-law, Son-in-law, Sister-in-law}
\label{ch:ka-art-aunt-daughter-in-law-son-in-law-sister-in-law}
\hlchaptermodality{226}

\begin{chapteropening}
\textbf{By the end of this chapter:} \emph{I can say art, aunt, daughter-in-law, son-in-law and sister-in-law.}

Complete the last lesson of chapter 226: I can say art, aunt, daughter-in-law, son-in-law and sister-in-law.
\end{chapteropening}

\section[\texorpdfstring{\kn{ಕಲೆ} (kale)}{ಕಲೆ (kale)}]{\kn{ಕಲೆ} (kale) — art}

\label{lesson:KA-C226-kale}



\begin{quote}
Before the new one: say the Kannada for sugarcane, then the Kannada for a firecracker.
\end{quote}

\subsection*{You'll want to know: \kn{ಕಲೆ}}
\textbf{\kn{ಕಲೆ}} — \emph{kale} — \textquotedblleft{}art\textquotedblright{}.

\begin{grammarlens}[title={What you've built}]
One more word for this chapter.
\end{grammarlens}

\subsection*{Your turn}
\begin{itemize}
\raggedright
  \item \emph{Say it:} \emph{kale}
  \item \emph{Say it:} \emph{kale}, once more
  \item \emph{Say it:} say \emph{kale}
  \item \emph{Recall:} say the Kannada for sugarcane, then the Kannada for a firecracker, then say \emph{kale} again
\end{itemize}

\begin{tcolorbox}[breakable,colback=teal!4,colframe=teal!35!black,title={Before you move on}]
What does \kn{ಕಲೆ} mean? (\textquotedblleft{}Art\textquotedblright{}.) Say it once more.
\end{tcolorbox}
\section[\texorpdfstring{\kn{ಅತ್ತೆ} (atte)}{ಅತ್ತೆ (atte)}]{\kn{ಅತ್ತೆ} (atte) — aunt (father's sister); mother-in-law}

\label{lesson:KA-C226-atte}



\begin{quote}
Before the new one: say the Kannada for a firecracker, then the Kannada for art.
\end{quote}

\subsection*{You'll want to know: \kn{ಅತ್ತೆ}}
\textbf{\kn{ಅತ್ತೆ}} — \emph{atte} — \textquotedblleft{}aunt (father's sister); mother-in-law\textquotedblright{}.

\begin{grammarlens}[title={What you've built}]
One more word for this chapter.
\end{grammarlens}

\subsection*{Your turn}
\begin{itemize}
\raggedright
  \item \emph{Say it:} \emph{atte}
  \item \emph{Say it:} \emph{atte}, once more
  \item \emph{Say it:} say \emph{atte}
  \item \emph{Recall:} say the Kannada for a firecracker, then the Kannada for art, then say \emph{atte} again
\end{itemize}

\begin{tcolorbox}[breakable,colback=teal!4,colframe=teal!35!black,title={Before you move on}]
What does \kn{ಅತ್ತೆ} mean? (\textquotedblleft{}Aunt (father's sister); mother-in-law\textquotedblright{}.) Say it once more.
\end{tcolorbox}
\section[\texorpdfstring{\kn{ಸೊಸೆ} (sose)}{ಸೊಸೆ (sose)}]{\kn{ಸೊಸೆ} (sose) — daughter-in-law}

\label{lesson:KA-C226-sose}



\begin{quote}
Before the new one: say the Kannada for art, then the Kannada for aunt.
\end{quote}

\subsection*{You'll want to know: \kn{ಸೊಸೆ}}
\textbf{\kn{ಸೊಸೆ}} — \emph{sose} — \textquotedblleft{}daughter-in-law\textquotedblright{}.

\begin{grammarlens}[title={What you've built}]
One more word for this chapter.
\end{grammarlens}

\subsection*{Your turn}
\begin{itemize}
\raggedright
  \item \emph{Say it:} \emph{sose}
  \item \emph{Say it:} \emph{sose}, once more
  \item \emph{Say it:} say \emph{sose}
  \item \emph{Recall:} say the Kannada for art, then the Kannada for aunt, then say \emph{sose} again
\end{itemize}

\begin{tcolorbox}[breakable,colback=teal!4,colframe=teal!35!black,title={Before you move on}]
What does \kn{ಸೊಸೆ} mean? (\textquotedblleft{}Daughter-in-law\textquotedblright{}.) Say it once more.
\end{tcolorbox}
\section[\texorpdfstring{\kn{ಅಳಿಯ} (aḷiya)}{ಅಳಿಯ (aḷiya)}]{\kn{ಅಳಿಯ} (aḷiya) — son-in-law}

\label{lesson:KA-C226-aliya}



\begin{quote}
Before the new one: say the Kannada for aunt, then the Kannada for daughter-in-law.
\end{quote}

\subsection*{You'll want to know: \kn{ಅಳಿಯ}}
\textbf{\kn{ಅಳಿಯ}} — \emph{aḷiya} — \textquotedblleft{}son-in-law\textquotedblright{}.

\begin{grammarlens}[title={What you've built}]
One more word for this chapter.
\end{grammarlens}

\subsection*{Your turn}
\begin{itemize}
\raggedright
  \item \emph{Say it:} \emph{aḷiya}
  \item \emph{Say it:} \emph{aḷiya}, once more
  \item \emph{Say it:} say \emph{aḷiya}
  \item \emph{Recall:} say the Kannada for aunt, then the Kannada for daughter-in-law, then say \emph{aḷiya} again
\end{itemize}

\begin{tcolorbox}[breakable,colback=teal!4,colframe=teal!35!black,title={Before you move on}]
What does \kn{ಅಳಿಯ} mean? (\textquotedblleft{}Son-in-law\textquotedblright{}.) Say it once more.
\end{tcolorbox}
\section[\texorpdfstring{\kn{ಅತ್ತಿಗೆ} (attige)}{ಅತ್ತಿಗೆ (attige)}]{\kn{ಅತ್ತಿಗೆ} (attige) — sister-in-law (elder brother's wife)}

\label{lesson:KA-C226-attige}



\begin{quote}
Before the new one: say the Kannada for daughter-in-law, then the Kannada for son-in-law.
\end{quote}

\subsection*{You'll want to know: \kn{ಅತ್ತಿಗೆ}}
\textbf{\kn{ಅತ್ತಿಗೆ}} — \emph{attige} — \textquotedblleft{}sister-in-law (elder brother's wife)\textquotedblright{}.

\begin{grammarlens}[title={What you've built}]
One more word for this chapter.
\end{grammarlens}

\subsection*{Your turn}
\begin{itemize}
\raggedright
  \item \emph{Say it:} \emph{attige}
  \item \emph{Say it:} \emph{attige}, once more
  \item \emph{Say it:} say \emph{attige}
  \item \emph{Recall:} say the Kannada for daughter-in-law, then the Kannada for son-in-law, then say \emph{attige} again
\end{itemize}

\begin{tcolorbox}[breakable,colback=teal!4,colframe=teal!35!black,title={Before you move on}]
What does \kn{ಅತ್ತಿಗೆ} mean? (\textquotedblleft{}Sister-in-law (elder brother's wife)\textquotedblright{}.) Say it once more.
\end{tcolorbox}
